\subsection{A criterion for formally smooth local algebras}

PROPOSITION 7. Let \(k_{0}\) be a ring, \(k\) a \(k_{0}\)-algebra, A a \(k\)-algebra, and \(\mathfrak{m}\) a maximal ideal of A. We assume that \(k\) and A/\(\mathfrak{m}\) are formally smooth over \(k_{0}\). For A to be formally smooth over \(k\) for the \(\mathfrak{m}\)-adic topology, it is necessary and sufficient that the following two conditions be satisfied:

\begin{enumerate}
\def\labelenumi{(\roman{enumi})}
\item
  the canonical homomorphism \(\mathsf{S}_{\mathrm{A} / \mathfrak{m}}(\mathfrak{m} / \mathfrak{m}^{2})\to \mathrm{gr}_{\mathfrak{m}}(\mathrm{A})\) is bijective;
\item
  the A/m-linear map
\end{enumerate}

\[
\omega : \mathrm{A} / \mathfrak {m} \otimes_ {k} \Omega_ {k _ {0}} (k) \longrightarrow \mathrm{A} / \mathfrak {m} \otimes_ {\mathrm{A}} \Omega_ {k _ {0}} (\mathrm{A})
\]

deduced from the canonical map \(k \to A\) is injective.

Let \(d_{k}:k\to \Omega_{k_{0}}(k)\) and \(d_{\mathrm{A}}:\mathrm{A}\to \Omega_{k_0}(\mathrm{A})\) be the universal \(k_{0}\)-derivations.

Suppose first that A is formally smooth over \(k\) for the \(\mathfrak{m}\)-adic topology. Then A is formally smooth over \(k_{0}\) for the \(\mathfrak{m}\)-adic topology (\(n^{\circ} 2\), prop. 3, a)), which is equivalent to (i) (\(n^{\circ} 5\), th. 2). Moreover, the \(k_{0}\)-derivation \(\lambda \mapsto 1 \otimes d_{k}(\lambda)\) of \(k\) in \(A/\mathfrak{m} \otimes_{k} \Omega_{k_{0}}(k)\) can be extended to a \(k_{0}\)-derivation of A into \(A/\mathfrak{m} \otimes_{k} \Omega_{k_{0}}(k)\) (\(n^{\circ} 2\), remark 2). There therefore exists an A-linear map \(u: \Omega_{k_{0}}(A) \to A/\mathfrak{m} \otimes_{k} \Omega_{k_{0}}(k)\) such that \(u(d_{A}(\lambda 1_{A})) = 1 \otimes d_{k}(\lambda)\) for every \(\lambda \in k\). The \(A/\mathfrak{m}\)-linear map \(A/\mathfrak{m} \otimes_{A} \Omega_{k_{0}}(A) \longrightarrow A/\mathfrak{m} \otimes_{k} \Omega_{k_{0}}(k)\) deduced from \(u\) is a retraction of \(\omega\), which proves (ii).

Conversely, suppose conditions (i) and (ii) are satisfied. Then A is formally smooth over \(k_{0}\) for the \(\mathfrak{m}\)-adic topology (\(n^{\circ}5\), th. 2) and the A-module \(\Omega_{k_{0}}(A)\) is projective (\(n^{\circ}5\), cor. 3 of th. 2). Fix an integer \(r \geqslant 0\) and consider the map, linear over \(A/\mathfrak{m}^{r}\),

\[
\omega_ {r}: \mathrm{A} / \mathfrak {m} ^ {r} \otimes_ {k} \Omega_ {k _ {0}} (k) \longrightarrow \mathrm{A} / \mathfrak {m} ^ {r} \otimes_ {\mathrm{A}} \Omega_ {k _ {0}} (\mathrm{A})
\]

deduced from the canonical map \(k\to\mathrm{A}\). Let \((\lambda_{i})_{i\in\mathrm{I}}\) be a family of elements of \(k\) such that the \(d_{k}(\lambda_{i})\) form a basis of the \(k\)-vector space \(\Omega_{k_{0}}(k)\). By (ii), the elements \(1 \otimes d_{A}(\lambda_{i}1_{A})\) are linearly independent in A/m \(\otimes_{\mathrm{A}} \Omega_{k_0}(\mathrm{A})\). By II, § 3, no. 2, cor. 1 and 2 of prop. 5, the \(1 \otimes d_{A}(\lambda_{i}1_{A})\) form a basis of a direct summand of the A/m \(^{r}\)-module A/m \(^{r} \otimes_{\mathrm{A}} \Omega_{k_0}(\mathrm{A})\). There therefore exists a map, linear over A/m \(^{r}\),

\[
u _ {r}: \mathrm{A} / \mathfrak {m} ^ {r} \otimes_ {\mathrm{A}} \Omega_ {k _ {0}} (\mathrm{A}) \longrightarrow \mathrm{A} / \mathfrak {m} ^ {r} \otimes_ {k} \Omega_ {k _ {0}} (k)
\]

such that \(u_r(1 \otimes d_A(\lambda_i 1_A)) = 1 \otimes d_k(\lambda_i)\) for every \(i\) , hence \(u_r \circ \omega_r = \mathrm{Id}\) .

Let us now verify that A is formally smooth over \(k\) for the m-adic topology. Let C be a \(k\)-algebra, N an ideal of square zero in C, and \(\pi: C \to C/N\) the canonical surjection; endow C and C/N with the discrete topology. Let \(\varphi: A \to C/N\) be a continuous homomorphism of \(k\)-algebras. Since A is formally smooth over \(k_0\) for the m-adic topology, there exists a continuous homomorphism of \(k_0\)-algebras \(\tilde{\varphi}_0: A \to C\) such that \(\pi \circ \tilde{\varphi}_0 = \varphi\). According to Proposition 1 of No. 1, the homomorphisms of \(k_0\)-algebras \(\tilde{\varphi}: A \to C\) such that \(\pi \circ \tilde{\varphi} = \varphi\) are the maps \(x \mapsto v(d_A(x)) + \tilde{\varphi}_0(x)\), where \(v\) ranges over \(\mathrm{Hom}_{A}(\Omega_{k_0}(A), N)\). It remains to choose \(v\) so that \(\tilde{\varphi}\) is a homomorphism of \(k\)-algebras. The map \(\lambda \mapsto \lambda 1_C - \tilde{\varphi}_0(\lambda 1_A)\) is a \(k_0\)-derivation from \(k\) into N (loc. cit.), hence can be written \(h \circ d_k\) with \(h \in \mathrm{Hom}_k(\Omega_{k_0}(k), N)\).

Let us choose an integer \(r\) such that the kernel of \(\varphi\) contains \(\mathfrak{m}^r\). The A-module N is annihilated by \(\mathfrak{m}^r\), and it suffices to take for \(v\) the composite of the sequence of homomorphisms

\[
\Omega_ {k _ {0}} (\mathrm{A}) \longrightarrow \mathrm{A} / \mathfrak {m} ^ {r} \otimes_ {\mathrm{A}} \Omega_ {k _ {0}} (\mathrm{A}) \xrightarrow {u _ {r}} \mathrm{A} / \mathfrak {m} ^ {r} \otimes_ {k} \Omega_ {k _ {0}} (k) \xrightarrow {h ^ {\prime}} \mathrm{N},
\]

where \(h'\) is deduced from \(h\). Indeed, we have, for \(\lambda \in k\):

\[
v \left(d _ {\mathrm{A}} \left(\lambda 1 _ {\mathrm{A}}\right)\right) = h ^ {\prime} u _ {r} \left(1 \otimes d _ {\mathrm{A}} \left(\lambda 1 _ {\mathrm{A}}\right)\right) = h ^ {\prime} \left(1 \otimes d _ {k} (\lambda)\right) = h \left(d _ {k} (\lambda)\right) = \lambda 1 _ {\mathrm{C}} - \tilde {\varphi} _ {0} \left(\lambda 1 _ {\mathrm{A}}\right).
\]

Remark 1.—When A is Noetherian, condition (i) means that the local ring \(A_{\mathfrak{m}}\) is regular (VIII, § 5, no. 2, th. 1).

PROPOSITION 8. Let \(k\) be a field and A a Noetherian local \(k\)-algebra. The following conditions are equivalent:

\begin{enumerate}
\def\labelenumi{(\roman{enumi})}
\item
  A is formally smooth over \(k\) for the \(\mathfrak{m}_{A}\)-adic topology;
\item
  A is regular and the \(\kappa_{\mathrm{A}}\)-linear map
\end{enumerate}

\[
\omega : \kappa_ {\mathrm{A}} \otimes_ {k} \Omega (k) \longrightarrow \kappa_ {\mathrm{A}} \otimes_ {\mathrm{A}} \Omega (\mathrm{A})
\]

deduced from the canonical injection k → A is injective ;

\begin{enumerate}
\def\labelenumi{(\roman{enumi})}
\setcounter{enumi}{2}
\item
  A is absolutely regular;
\item
  for every radical extension \(k'\) of \(k\) , of finite degree and height \(\leqslant 1\) , the local ring \(A_{(k')}\) is regular.
\item
  (ii) \(\Leftrightarrow\) (i): it suffices to apply Proposition 7 and Remark 1 above, taking \(k_{0}\) to be the prime subfield of \(k\); indeed, \(k\) and \(\kappa_{\mathrm{A}}\) are formally smooth over \(k_{0}\) (\(n^{\circ}3\), cf. th. 1).
\end{enumerate}

(i)⇒(iii): this follows from cor. 2 of th. 2 (no. 5).

(iii)⇒(iv): this follows from def. 1 of § 6, no. 4.

If \(k\) is of characteristic 0, it follows from Corollary 1 of Theorem 2 (No. 5) that (iv) implies (i), whence the proposition in this case. Suppose \(k\) has characteristic \(p \neq 0\) and prove (iv)⇒(ii). Let \(k'\) be a radical extension of \(k\) of finite degree and height \(\leqslant 1\). If A and \(\mathrm{A}_{(k')}\) are regular, the canonical map \(\kappa_{\mathrm{A}} \otimes_{k'^{p}} \Omega_{k^{p}}(k'^{p}) \longrightarrow \kappa_{\mathrm{A}} \otimes_{\mathrm{A}} \Omega(\mathrm{A})\) is injective (No. 6, Proposition 6). By Theorem 1(b) of A, V, p.~97, applied to the extension \(k\) of \(k^{p}\), the \(k\)-vector space \(\Omega(k)\), which coincides with \(\Omega_{k^{p}}(k)\), is the increasing filtered union of the subspaces \(k \otimes_{k'^{p}} \Omega_{k^{p}}(k'^{p})\), where \(k'\) ranges over the finite radical extensions of \(k\) of height \(\leqslant 1\) contained in a fixed algebraic closure of \(k\). The assertion (ii) follows from this.

Remark 2.—Let \(k\) be a field and A a \(k\)-algebra such that the ring \(\mathrm{A}_{(k')}\) is regular for every radical extension \(k'\) of \(k\), of finite degree and height \(\leqslant 1\); then A is absolutely regular. Indeed, let \(k'\) be such an extension; for every maximal ideal \(\mathfrak{m}\) of A, the ring \(k'\otimes_k\mathrm{A}_{\mathfrak{m}}\) is identified with a ring of fractions of \(\mathrm{A}_{(k')}\), hence is regular. By Proposition 8 above and Proposition 6 of § 6, No. 4, the algebra A is absolutely regular.

